
\documentclass{article}
\usepackage{colortbl}
\usepackage{makecell}
\usepackage{multirow}
\usepackage{supertabular}

\begin{document}

\newcounter{utterance}

\centering \large Interaction Transcript for game `guesswhat', experiment `Abs\_Level\_2', episode 5 with reka{-}flash{-}3.1{-}t0.0.
\vspace{24pt}

{ \footnotesize  \setcounter{utterance}{1}
\setlength{\tabcolsep}{0pt}
\begin{supertabular}{c@{$\;$}|p{.15\linewidth}@{}p{.15\linewidth}p{.15\linewidth}p{.15\linewidth}p{.15\linewidth}p{.15\linewidth}}
    \# & $\;$A & \multicolumn{4}{c}{Game Master} & $\;\:$B\\
    \hline

    \theutterance \stepcounter{utterance}  
    & & \multicolumn{4}{p{0.6\linewidth}}{
        \cellcolor[rgb]{0.9,0.9,0.9}{
            \makecell[{{p{\linewidth}}}]{
                \texttt{\tiny{[P1$\langle$GM]}}
                \texttt{Let's play a guessing game. Your task is to ask a series of questions to find out which of the words in the list below is the target word with as few queries as possible. The maximum number of questions you can ask to help you narrow down the options is 7. Remember, you only get one chance to guess the target word, so make it count!} \\
\\ 
\texttt{The rules of the game are: } \\
\\ 
\texttt{1. Only ask one question per turn. DO NOT ASK MORE THAN ONE QUESTION.} \\
\\ 
\texttt{2. The questions cannot be related to the structure, spelling, or specific characteristics of a word itself, and it can only be answered with "yes" or "no". } \\
\\ 
\texttt{3. Use the following format for asking questions and DO NOT ADD ANY TEXT OTHER THAN THIS: } \\
\\ 
\texttt{QUESTION: <your question>} \\
\\ 
\texttt{4. You have only one opportunity to guess the target word. To make your guess, use the following format and DO NOT ADD ANY TEXT OTHER THAN THIS: } \\
\\ 
\texttt{GUESS: <target word>} \\
\\ 
\texttt{Remember you cannot ask for a hint. Make your question brief. Let’s start!} \\
\\ 
\texttt{List of candidate words: {[}'sixth', 'fourth', 'third', 'tenth', 'name', 'word', 'verb', 'adverb'{]}} \\
            }
        }
    }
    & \\ \\

    \theutterance \stepcounter{utterance}  
    & \multicolumn{4}{p{0.6\linewidth}}{
        \cellcolor[rgb]{0.9,0.9,0.9}{
            \makecell[{{p{\linewidth}}}]{
                \texttt{\tiny{[P1$\rangle$GM]}}
                \texttt{<reasoning>} \\
\texttt{Let's see. I need to figure out which word from the list is the target using up to 7 yes/no questions. The list is {[}'sixth', 'fourth', 'third', 'tenth', 'name', 'word','verb', 'adverb'{]}.} \\
\\ 
\texttt{First, I should look for common categories or characteristics. The words can be divided into ordinal numbers (sixth, fourth, third, tenth) and other words (name, word, verb, adverb). Maybe start by separating these two groups. If I can determine whether the target is an ordinal or not, that would cut the list in half.} \\
\\ 
\texttt{So, first question: Is the target word an ordinal number? That's a yes/no question. Let's go with that. If the answer is yes, then the possible words are sixth, fourth, third, tenth. If no, then name, word, verb, adverb. } \\
\\ 
\texttt{Assuming the answer is yes (since there are four ordinals and four others, but maybe the target is more likely to be in one group? Not sure. But since the user hasn't given any hints, I have to proceed with equal probability. Wait, but the user hasn't specified any bias. So maybe the first question is best to split the list into two equal parts. Since there are 8 words, splitting into 4 and 4. So the first question should aim to do that. } \\
\\ 
\texttt{Alternatively, maybe the ordinals are more likely? Not necessarily. Let's proceed with the first approach. } \\
\\ 
\texttt{So first question: Is the target word an ordinal number? Let's ask that. } \\
\\ 
\texttt{If the answer is yes, then next step: among the four ordinals, how to narrow down. Maybe ask about the position. For example, is it the third position? If yes, then third. If no, then maybe ask if it's the fourth. But that might take more questions. Alternatively, ask if it's a single{-}syllable word. Wait, but the rules say questions can't be about structure or spelling. So can't ask about syllables. Hmm. } \\
\\ 
\texttt{Alternatively, think of other characteristics. For example, ordinals can be followed by 'th' except for first, second, third, fourth, fifth, etc. Wait, but sixth, fourth, third, tenth all end with 'th' except maybe tenth? Wait, tenth ends with 'th' as well.} \\
            }
        }
    }
    & & \\ \\

    \theutterance \stepcounter{utterance}  
    & & & \multicolumn{2}{p{0.3\linewidth}}{
        \cellcolor[rgb]{0.9,0.9,0.9}{
            \makecell[{{p{\linewidth}}}]{
                \texttt{\tiny{[GM$|$GM]}}
                \texttt{Invalid format. Guesser must use the form 'QUESTION: ' or 'GUESS: '.} \\
            }
        }
    }
    & & \\ \\

    \theutterance \stepcounter{utterance}  
    & & & \multicolumn{2}{p{0.3\linewidth}}{
        \cellcolor[rgb]{0.9,0.9,0.9}{
            \makecell[{{p{\linewidth}}}]{
                \texttt{\tiny{[GM$|$GM]}}
                \texttt{abort game} \\
            }
        }
    }
    & & \\ \\

\end{supertabular}
}

\end{document}
